% Source: https://artofproblemsolving.com/community/c3h1919033_no_test_points
% Stable identity: aops-topic-1919033-q01
(a) A factor $x-r$ is negative for $x<r$ and positive for $x>r$. For $x>r_k$ every factor is positive, so $p(x)$ has the sign of $c$. On an open interval containing no root, no factor changes sign, so $p$ keeps its sign. When $x$ moves from just above $r_i$ to just below it, only the factor $(x-r_i)^{m_i}$ changes: $x-r_i$ changes from positive to negative, so $(x-r_i)^{m_i}$ changes sign exactly when $m_i$ is odd. So we can write down all the signs by starting with the sign of $c$ on the far right and moving left, switching at each root of odd multiplicity and keeping the sign at each root of even multiplicity.

(b) Since $1-2-5+6=0$, $x=1$ is a root, and
\[
x^3-2x^2-5x+6=(x-1)(x^2-x-6)=(x+2)(x-1)(x-3).
\]
The roots $-2,1,3$ are simple and $c=1>0$, so the signs from right to left are $+$ on $(3,\infty)$, $-$ on $(1,3)$, $+$ on $(-2,1)$ and $-$ on $(-\infty,-2)$. Hence
\[
x^3-2x^2-5x+6>0\iff \boxed{-2<x<1\ \text{or}\ x>3}.
\]
For $(x+2)(x-1)^2(x-3)$ the leading coefficient is $1$, the roots $-2$ and $3$ are simple and $1$ is a double root. From right to left the signs are $+$ on $(3,\infty)$, $-$ on $(1,3)$, still $-$ on $(-2,1)$ (no change at the double root) and $+$ on $(-\infty,-2)$. Hence
\[
(x+2)(x-1)^2(x-3)<0\iff \boxed{-2<x<3\ \text{and}\ x\ne1}.
\]
